If $A$ is a valuation ring, any two nonzero elements $a,b$ satisfy $a/b\in A$ or $b/a\in A$, so their principal ideals are comparable. If $\mathfrak{a}\nsubseteq\mathfrak{b}$, choose $a\in\mathfrak{a}\setminus\mathfrak{b}$. For every $b\in\mathfrak{b}$, comparability forces $(b)\subseteq(a)$, since the reverse inclusion would put $a$ in $\mathfrak{b}$. Hence $\mathfrak{b}\subseteq\mathfrak{a}$.

Conversely, for $0\neq a/b\in K$, comparability of $(a)$ and $(b)$ gives $a/b\in A$ or $b/a\in A$. Thus $A$ is a valuation ring.

Ideals of $A/\mathfrak{p}$ correspond to ideals of $A$ containing $\mathfrak{p}$, so they remain totally ordered. The domain $A/\mathfrak{p}$ is therefore a valuation ring of its fraction field. Also $A\subseteq A_{\mathfrak{p}}\subseteq K$, so every nonzero element of $K$ or its inverse belongs to $A_{\mathfrak{p}}$.
